% =====================================================================
\section{Справочник}
% =====================================================================
\subsection*{Выводы, с которыми нужно быть осторожным}

\begin{center}\small
\begin{tabularx}{\textwidth}{@{}l X@{}}
\toprule
Выводы & Особенность \\
\midrule
GPIO0, 3, 45, 46 & «загрузочные»: влияют на запуск; GPIO0 --- кнопка BOOT \\
GPIO19, 20 & USB (D$-$, D$+$) \\
GPIO26--32 & память программ --- не использовать \\
GPIO35--37 & заняты памятью PSRAM на модулях с Octal PSRAM (N8R8, N16R8) \\
GPIO43, 44 & UART0 (TX, RX) --- сообщения при загрузке \\
GPIO11--20 & ADC2 --- не работает при включённом Wi-Fi \\
GPIO22--25 & не существуют \\
\bottomrule
\end{tabularx}
\end{center}

\subsection*{Шпаргалка по языку C++}

\begin{center}\small
\begin{tabularx}{\textwidth}{@{}l X@{}}
\toprule
Запись & Смысл\\
\midrule
\code{int x = 5;} & переменная целого типа со значением 5\\
\code{const int PIN = 4;} & константа: значение не меняется\\
\code{x = x + 1;}\quad \code{x++;}\quad \code{x += 1;} & три способа увеличить \code{x} на 1\\
\code{a / b}, \code{a \% b} & деление (для целых --- нацело) и остаток\\
\code{==  !=  <  >  <=  >=} & сравнение\\
\code{\&\&  ||  !} & логические «и», «или», «не»\\
\code{if (усл) \{ \ldots \} else \{ \ldots \}} & выполнить одну из двух веток\\
\code{for (int i = 0; i < N; i++) \{ \ldots \}} & повторить $N$ раз, \code{i} = 0, 1, \ldots, $N-1$\\
\code{while (усл) \{ \ldots \}} & повторять, пока условие верно; \code{break} --- выйти\\
\code{int a[] = \{1, 2, 3\};}\quad \code{a[0]} & массив и его первый элемент\\
\code{int sum(int a, int b) \{ return a + b; \}} & функция с параметрами и результатом\\
\code{void hello() \{ \ldots \}} & функция без результата\\
\code{// текст} & комментарий\\
\bottomrule
\end{tabularx}
\end{center}

\subsection*{Шпаргалка по функциям Arduino-ESP32}

\begin{center}\small
\begin{tabularx}{\textwidth}{@{}l X l@{}}
\toprule
Функция & Что делает & Урок\\
\midrule
\code{pinMode(pin, OUTPUT / INPUT\_PULLUP)} & настроить вывод & \ref{les:led}, \ref{les:gpio-in}\\
\code{digitalWrite(pin, HIGH / LOW)} & выдать 1 или 0 & \ref{les:led}\\
\code{digitalRead(pin)} & прочитать 0 или 1 & \ref{les:gpio-in}\\
\code{millis()}, \code{delay(ms)} & время с момента включения; пауза & \ref{les:led}\\
\code{ledcAttach(pin, freq, bits)} & включить ШИМ на выводе & \ref{les:pwm}\\
\code{ledcWrite(pin, duty)} & задать коэффициент заполнения & \ref{les:pwm}\\
\code{ledcWriteTone(pin, freq)} & звук заданной частоты & \ref{les:pwm}\\
\code{WiFi.begin(ssid, pass)} & подключиться к сети & \ref{les:radio}\\
\code{WiFi.softAP(ssid, pass)} & создать свою сеть & \ref{les:radio}\\
\code{MDNS.begin("имя")} & имя \code{имя.local} & \ref{les:radio}\\
\code{server.on(путь, функция)} & обработчик адреса веб-сервера & \ref{les:web}\\
\code{Serial.begin()}, \code{print()}, \code{readStringUntil()} & связь с компьютером & \ref{les:serial}\\
\code{attachInterrupt(pin, isr, FALLING)} & прерывание от вывода & \ref{les:gpio-in}\\
\code{analogReadMilliVolts(pin)} & измерить напряжение, мВ & \ref{les:adc}\\
\code{Wire.begin(sda, scl)} & запустить шину I\textsuperscript{2}C & \ref{les:i2c}\\
\code{xTaskCreatePinnedToCore(...)} & создать задачу FreeRTOS & \ref{les:rtos}\\
\code{esp\_deep\_sleep\_start()} & уйти в глубокий сон & \ref{les:sleep}\\
\bottomrule
\end{tabularx}
\end{center}

\subsection*{Полезные формулы}

\begin{center}\small
\renewcommand{\arraystretch}{1.6}
\begin{tabular}{@{}ll@{}}
\toprule
Закон Ома & $I = U/R$, \quad $R = U/I$\\
Резистор для светодиода & $R = (U_\text{пит} - U_F)/I$\\
Частота и период & $f = 1/T$\\
Коэффициент заполнения ШИМ & $D = t_\text{вкл}/T$, \quad $U_\text{ср} = D\cdot U_\text{пит}$\\
Число вариантов в $N$ битах & $2^N$\\
Шаг квантования АЦП & $q = U_\text{макс}/2^N$\\
Длина волны & $\lambda = c/f$\\
Делитель напряжения & $U_\text{вых} = U_\text{вх}\cdot R_2/(R_1+R_2)$\\
\bottomrule
\end{tabular}
\end{center}

\subsection*{Что изучать дальше}

\begin{figure}[H]
\centering
\begin{tikzpicture}[node distance=5mm and 6mm, b/.style={blk,minimum width=28mm}]
  \node[b,fill=blkcpu] (base) {Этот курс\\ (уроки 1--14)};
  \node[b,fill=blkrf,above right=4mm and 12mm of base] (ble) {Bluetooth LE\\ \scriptsize связь с телефоном};
  \node[b,fill=blkrf,right=12mm of base] (mqtt) {MQTT, умный дом,\\ обновление по Wi-Fi};
  \node[b,fill=blkrf,below right=4mm and 12mm of base] (usb) {USB-клавиатура,\\ MIDI-инструмент};
  \node[b,fill=blkper,right=of ble] (cam) {Камера,\\ цветной дисплей};
  \node[b,fill=blkper,right=of mqtt] (ai) {Нейросети\\ на микроконтроллере};
  \node[b,fill=blkper,right=of usb] (idf) {ESP-IDF ---\\ «профессиональный» уровень};
  \foreach \t in {ble,mqtt,usb} \draw[arr] (base.east) -- (\t.west);
  \draw[arr] (ble) -- (cam); \draw[arr] (mqtt) -- (ai); \draw[arr] (usb) -- (idf);
\end{tikzpicture}
\caption{Куда двигаться после этого курса}
\end{figure}

\subsection*{Полезные источники}
\begin{itemize}
  \item Документация на ESP32-S3 --- \url{https://www.espressif.com/en/support/documents/technical-documents}
  \item Документация Arduino-ESP32 --- \url{https://docs.espressif.com/projects/arduino-esp32/}
  \item Руководство ESP-IDF --- \url{https://docs.espressif.com/projects/esp-idf/en/latest/esp32s3/}
  \item Описание платы --- \emph{ESP32-S3-DevKitC-1 User Guide} на \url{https://docs.espressif.com}
  \item Учебник по HTML и CSS --- \url{https://developer.mozilla.org/ru/docs/Learn}
\end{itemize}
